\input{./._relpath-to-latexroot.ltx}
\documentclass[\latexroot/\projectname]{subfiles}
\input{\latexroot/@resources/subfile-setup.ltx}

\begin{document}

\FloatBarrier
\section{Comparing fiscal stimulus policies}
\whenintegrated{\label{sec:comparing}} % integrated doc: SST for labels


This section analyzes the macroeconomic effects of MP, CBI, and RN shocks extracted from target-rate surprises. We report impulse responses under two identification approaches. The first corresponds to the \emph{standard HFI} approach, which treats the target-rate surprise as an exogenous MP shock and does not separately identify RN and CBI components. The second relies on the sign- and magnitude-restriction identification strategy, which disentangles MP, RN, and CBI shocks from target-rate surprises. MP shocks identified under both approaches are normalized to correspond to a one percentage point increase in the one-year government bond yield, while CBI and RN shocks are scaled to match the size of the normalized pure MP shock in standard deviation units.\footnote{Standardizing the size of MP shocks facilitates comparability across identification approaches, as well as with the existing literature, where a one percentage point increase in interest rates is commonly used as a benchmark. Scaling CBI and RN shocks to match the size of the normalized pure MP shock is particularly relevant for FG surprises, as RN shocks tend to generate much smaller movements in the one-year government bond yield. Normalizing such shocks to a one percentage point increase would require disproportionately large innovations, which would distort the scale of the impulse responses and hinder their visual comparability. We report impulse responses without standardization in \hyperlink{Figure B.23}{Figure~B.23} of \hyperlink{Appendix B}{Appendix~B}.} The corresponding impact responses of the high-frequency surprises—by construction consistent with the imposed sign and magnitude restrictions—are reported in \hyperlink{Table 3}{Table~3}.

\begin{table}[t]
\hypertarget{Panel AT3}{}
\hypertarget{Panel BT3}{}
\hypertarget{Table 3}{}
\begin{center}
\resizebox{16cm}{!} {
\begin{tabular}{c c c c c c c c c}
\toprule
& \multicolumn{2}{c}{Panel A. Standard HFI}
& \multicolumn{6}{c}{Panel B. Sign- and Magnitude-Restriction Identification} \\
\cmidrule(lr){2-3}\cmidrule(lr){4-9}
High Frequency
& \multicolumn{2}{c}{Target-Rate Shock}
& \multicolumn{2}{c}{Pure Target-Rate Shock}
& \multicolumn{2}{c}{RN Shock}
& \multicolumn{2}{c}{CBI Shock} \\
\cmidrule(lr){2-3}\cmidrule(lr){4-5}\cmidrule(lr){6-7}\cmidrule(lr){8-9}
Surprise
& Mean & ($16^{pct}$, $84^{pct}$)
& Mean & ($16^{pct}$, $84^{pct}$)
& Mean & ($16^{pct}$, $84^{pct}$)
& Mean & ($16^{pct}$, $84^{pct}$) \\
\midrule
Target-Rate surprise
& 0.50 & (0.48, 0.52)
&  & 
&   & 
&   &  \\
RES
&  & 
& 0.68 & (0.47, 0.80)
& -0.04 & (-0.33, 0.26)
& 0.34 & (0.06, 0.65) \\
FIT
&   & 
& -0.04 & (-0.15, 0.07)
& 0.25 & (0.17, 0.31)
& 0.16 & (0.04, 0.27) \\
Stock market surprise
&   & 
& -4.68 & (-7.17, -1.74)
& -4.10 & (-6.64, -1.38)
& 5.31 & (2.47, 7.70) \\
\bottomrule
\end{tabular}}
\caption{Impact Responses of Target-Rate and Stock Market Surprises}
\end{center}
\begin{footnotesize}
\textit{Notes:} Posterior means and 16th and 84th percentiles. Units are percentage points. RES and FIT denote the residual and fitted components of the target-rate surprise, respectively. Target-rate shocks identified under both approaches are normalized to correspond to a one percentage point increase in the one-year government bond yield. CBI and RN shocks are scaled to match the size of the normalized pure target-rate shock in standard deviation units.
\end{footnotesize}
\end{table}


\hyperlink{Panel BF1}{Panel A} of \hyperlink{Figure 1}{Figure 1} presents impulse responses to a contractionary target-rate shock \emph{without} disentangling the RN and CBI components. We report median responses over a 72-month horizon, together with 68\% highest posterior density intervals (HPDIs) corresponding to the 16th and 84th percentiles of the posterior distribution.\footnote{The use of 68\% confidence intervals is standard in related studies (see \hyperlink{Eichenbaum and Evans (1995)}{Eichenbaum and Evans, 1995}; \hyperlink{Sims and Zha (1999)}{Sims and Zha, 1999}; \hyperlink{Kim and Roubini (2000)}{Kim and Roubini, 2000}; \hyperlink{Kim (2001)}{Kim, 2001}; \hyperlink{Faust and Rogers (2003)}{Faust and Rogers, 2003}; \hyperlink{Scholl and Uhlig (2008)}{Scholl and Uhlig, 2008}; \hyperlink{Rogers et al. (2018)}{Rogers et al., 2018}; \hyperlink{Breitenlechner et al. (2021)}{Breitenlechner et al., 2021}).}   The results suggest that these shocks are initially associated with an expansion in output, rising prices, a depreciation of the domestic currency (i.e., an increase in the exchange rate), and higher inflation expectations. Although output and prices eventually decline, the responses appear to be relatively attenuated and reach their trough with a substantial lag. The peak contraction in output is about 0.27\% and occurs around two years after the shock, while the peak decline in prices is roughly 0.11\% and materializes more than four years after the disturbance. A similar pattern emerges for the Mexican peso, which exhibits a delayed and hump-shaped appreciation, reaching its maximum more than two years after the shock, with a magnitude of about 0.97\%.

These dynamics stand in contrast to the predictions of standard macroeconomic models and to the empirical evidence reported by \hyperlink{Carrillo et al. (2025)}{Carrillo et al. (2025)} for Mexico. In their estimates, a 1\% MP shock leads to output contractions of 0.30--0.65\% peaking within 12--15 months, inflation declines of 0.15--0.22\% within two years, and an immediate exchange rate appreciation of between 1--3.5\%. They also diverge from the international evidence documented by \hyperlink{Deb et al. (2023)}{Deb et al. (2023)}, who show that, across advanced and emerging economies, the maximum effects of a 1\% monetary tightening typically materialize within two years, with output falling by 0.20--0.40\% and inflation by around 0.20\% in EMEs.


\hyperlink{Panel BF1}{Panel~B} of \hyperlink{Figure~1}{Figure~1} presents impulse responses obtained after disentangling MP, RN, and CBI shocks. The first column reports impulse responses to a contractionary target-rate shock that is now purged of RN and CBI components. Notably, once these components are isolated, the previously documented output, price, and exchange-rate puzzles largely disappear. In addition, inflation expectations exhibit a negative, albeit muted, response, suggesting that they remain relatively anchored following the unexpected tightening. 

\begin{figure}[t]
\hypertarget{Figure 2}{}
\hypertarget{Panel AF2}{}
\hypertarget{Panel BF2}{}
\begin{center}
\resizebox{13.5cm}{!} {
\includegraphics[width=\linewidth]{Figures/Figure_1_Mex.pdf}
}
\caption[5pt]{Impulse Responses to MP, CBI, and RN shocks Disentangled from Target-Rate Surprises}
\end{center}
\begin{footnotesize} 
Notes: Responses  are  presented  for a 72-month horizon with the associated 68\% highest posterior density intervals (HPDIs). Target-rate shocks are normalized to imply a one percentage point increase in the one-year government bond yield. The CBI and RN shocks are scaled to match the magnitude of the normalized pure target-rate shock in standard deviation units.
\end{footnotesize} 
\end{figure}

Compared to the results in \hyperlink{Panel BF1}{Panel~A}, the contraction in output is substantially stronger and materializes earlier, with the peak decline occurring seven months sooner and reaching around 0.41\% below the baseline. A similar pattern emerges for prices, whose peak contraction is more than twice as pronounced (around 0.27\%) and occurs roughly two years after the disturbance—approximately half the time required in \hyperlink{Panel BF1}{Panel~A}. The exchange-rate response also changes markedly. Instead of the small and delayed appreciation observed following the raw disturbance, the Mexican peso now appreciates by approximately 4.25\% one month after the shock—a magnitude more than four times larger—before gradually returning toward its pre-shock level. Stock prices display a stronger impact response as well, declining by about 4.27\% compared to roughly 1.6\% after the unpurged shock. Finally, although both the one-year and ten-year government bond yields increase on impact, their responses appear less persistent than in \hyperlink{Panel BF1}{Panel~A}, suggesting that the tightening implied by the pure shock is shorter-lived and more consistent with a temporary policy adjustment. Overall, these estimates are more in line with the evidence reported by \hyperlink{Carrillo et al. (2025)}{Carrillo et al. (2025)} and \hyperlink{Deb et al. (2023)}{Deb et al. (2023)}, both in terms of magnitude and timing of the transmission mechanism.



We now describe the impulse responses to favorable CBI shocks reported in the third column of \hyperlink{Panel BF1}{Panel~B} of \hyperlink{Figure~1}{Figure~1}. In line with \hyperlink{Jarociński and Karadi (2025)}{Jarociński and Karadi (2020, 2025)}, favorable CBI shocks are followed by a persistent increase in the one-year government bond yield, an expansion in output, higher prices, and rising inflation expectations. Notably, the Mexican peso depreciates against the U.S.\ dollar in response to these shocks. Stock prices tend to rise on impact, while the ten-year government bond yield exhibits a positive, albeit muted, response. In particular, the depreciation of the domestic currency in response to favorable CBI shocks may be consistent with upward revisions in inflation expectations, as emphasized by \hyperlink{Gurkaynak et al. (2021)}{G\"urkaynak et al. (2021)} and \hyperlink{Stavrakeva and Tang (2019)}{Stavrakeva and Tang (2019)}.\footnote{An alternative explanation relates to Mexico’s tight integration with the United States. To the extent that the favorable information conveyed by the central bank is associated with positive spillovers from the U.S. economy—potentially accompanied by an appreciation of the U.S. dollar—then the estimated impulse responses may partly capture this external dollar-strength component, which would mechanically translate into a depreciation of the peso.}\textsuperscript{,}\footnote{While our findings are broadly consistent with U.S. evidence (\hyperlink{Pinchetti and Szczepaniak (2024)}{Pinchetti and Szczepaniak, 2024}; \hyperlink{Gründler et al. (2023)}{Gr\"undler et al., 2023}; \hyperlink{Stavrakeva and Tang (2019)}{Stavrakeva and Tang, 2019}), evidence from the euro area suggests a different pattern: favorable CBI shocks are associated with an appreciation of the domestic currency (\hyperlink{Kerssenfischer (2022)}{Kerssenfischer (2022)}).}  At the same time, the muted response of the ten-year government bond yield following a favorable CBI shock may reflect an offsetting mechanism. The stimulative effects of favorable information about economic fundamentals may reduce the term premium, thereby counteracting the upward pressure on long-term rates induced by the policy tightening.\footnote{The term premium is the additional compensation that investors require for holding longer-term instruments instead of rolling over short-term debt. According to \hyperlink{Carrillo et al. (2020)}{Carrillo et al. (2020)}, the response of the term premium depends on whether the information conveyed by the announcement is perceived as negative or positive for the economic outlook. The term premium tends to rise when news signals a deterioration in fundamentals and to fall when news points to a more favorable outlook.}

Taken together, these responses are consistent with an interpretation in which the central bank reveals favorable information about economic fundamentals while simultaneously tightening policy to lean against emerging overheating and inflationary pressures. Such pressures may operate through an aggregate demand channel and/or through exchange-rate pass-through effects. The persistent increase in the one-year government bond yield is consistent with a systematic policy response aimed at containing these pressures. While the tightening does not fully offset the initial expansionary effects of the news, it gradually neutralizes them over the medium term.

We now describe the impulse responses to hawkish RN shocks reported in the second column of \hyperlink{Panel BF1}{Panel~B} of \hyperlink{Figure~1}{Figure~1}. These shocks are followed by increases in interest rates, an expansion in output, higher prices, rising inflation expectations, and a depreciation of the domestic currency. These dynamics closely resemble those observed following favorable CBI shock and are consistent with the findings of \hyperlink{Bauer and Swanson (2023a, 2023b)}{Bauer and Swanson (2023a, 2023b)}. Overall, they may reflect a scenario in which the central bank tightens policy more aggressively than markets had anticipated in response to improving economic conditions or perceived upside risks to inflation. As in the CBI case, the tightening does not fully offset the expansionary effects associated with the underlying fundamentals but instead gradually neutralizes them over the medium term. 

Overall, these findings underscore the importance of disentangling MP shocks from their informational components. When target-rate surprises are treated as exogenous innovations, the resulting impulse responses may appear attenuated, delayed, or even puzzling, thereby obscuring the transmission mechanism of a conventional tightening. Once RN and CBI components are  isolated, the estimated effects of target-rate shocks become more consistent with standard theoretical predictions. In the next section, we apply the same decomposition framework to FG surprises, where distinguishing between policy and informational channels appears to be similarly important.




\subsection{Macroeconomic Effects of MP, RN, and CBI Shocks Disentangled from FG Surprises}

This section analyzes the macroeconomic effects of MP, CBI, and RN shocks extracted from FG surprises. As in the previous section, we present impulse responses under two identification approaches. The first corresponds to the \emph{standard HFI} approach, which treats the FG surprise as an exogenous MP shock and does not separately identify RN and CBI components. The second relies on a sign- and magnitude-restriction identification strategy that disentangles MP, RN, and CBI shocks from FG surprises. MP shocks identified under both approaches are normalized to correspond to a one percentage point increase in the one-year government bond yield. CBI and RN shocks are scaled to match the size of the normalized pure MP shock in standard deviation units.\footnote{\hyperlink{Figure B.24}{Figure~B.24} of \hyperlink{Appendix B}{Appendix~B} reports impulse responses without standardization.} The corresponding impact responses of the high-frequency surprises are reported in \hyperlink{Table 4}{Table~4}.

\begin{table}[t]
\hypertarget{Panel AT3}{}
\hypertarget{Panel BT3}{}
\hypertarget{Table 3}{}
\begin{center}
\resizebox{16cm}{!} {
\begin{tabular}{c c c c c c c c c}
\toprule
& \multicolumn{2}{c}{Panel A. Standard HFI}
& \multicolumn{6}{c}{Panel B. Sign- and Magnitude-Restriction Identification} \\
\cmidrule(lr){2-3}\cmidrule(lr){4-9}
High Frequency
& \multicolumn{2}{c}{FG Shock}
& \multicolumn{2}{c}{Pure FG Shock}
& \multicolumn{2}{c}{RN Shock}
& \multicolumn{2}{c}{CBI Shock} \\
\cmidrule(lr){2-3}\cmidrule(lr){4-5}\cmidrule(lr){6-7}\cmidrule(lr){8-9}
Surprise
& Mean & ($16^{pct}$, $84^{pct}$)
& Mean & ($16^{pct}$, $84^{pct}$)
& Mean & ($16^{pct}$, $84^{pct}$)
& Mean & ($16^{pct}$, $84^{pct}$) \\
\midrule
FG surprise
& 1.06 & (1.01, 1.10)
&  & 
&   & 
&   &  \\
RES
&  & 
& 1.47 & (1.02, 1.73)
& -0.24 & (-0.87, 0.46)
& 0.73 & (0.12, 1.41) \\
FIT
&   & 
& -0.12 & (-0.35, 0.14)
& 0.55 & (0.36, 0.65)
& 0.33 & (0.07, 0.58) \\
Stock market surprise
&   & 
& -4.40 & (-6.93, -1.57)
& -4.02 & (-6.51, -1.24)
& 5.69 & (2.95, 7.88) \\
\bottomrule
\end{tabular}}
\caption{Impact Responses of FG and Stock Market Surprises}
\end{center}
\begin{footnotesize}
\textit{Notes:} Posterior means and 16th and 84th percentiles. Units are percentage points. RES and FIT denote the residual and fitted components of the FG surprise, respectively. FG shocks identified under both approaches are normalized to correspond to a one percentage point increase in the one-year government bond yield. CBI and RN shocks are scaled to match the size of the normalized pure FG shock in standard deviation units.
\end{footnotesize}
\end{table}

\hyperlink{Panel BF1}{Panel~A} of \hyperlink{Figure~2}{Figure~2} displays in red impulse responses to a contractionary FG shock identified without disentangling RN and CBI components. For comparison, we also report in green the impulse responses to a contractionary pure target-rate shock, normalized to a one percentage point increase in the one-year government bond yield. The impulse responses to FG shocks do not appear to exhibit a clear attenuation bias potentially arising from RN and CBI components, with the possible exception of inflation expectations, which tend to rise in the short run. Compared with target-rate shocks, however, unpurged FG shocks are associated with a disproportionately large and more immediate contraction in domestic activity. Output declines by around 0.44\% one month after the shock, whereas the peak contraction following a pure target-rate shock (around 0.41\%) occurs only after about 19 months. Stock prices also appear to respond  more strongly following unpurged FG shocks. While these results could point to a more important role for FG in shaping macroeconomic outcomes, they may also reflect the amplification associated with the \emph{FG puzzle} discussed in the theoretical literature (see \hyperlink{Del Negro et al. (2023)}{Del Negro et al., 2023}), whereby revisions in expectations about the future policy path can generate macroeconomic effects that are disproportionately large relative to those of conventional MP shocks. At the same time, the relatively rapid response of output contrasts with the conventional view that monetary policy affects real activity with substantial lags.

\begin{figure}[t]
\hypertarget{Figure 2}{}
\hypertarget{Panel AF2}{}
\hypertarget{Panel BF2}{}
\begin{center}
\resizebox{13.5cm}{!} {
\includegraphics[width=\linewidth]{Figures/Figure_2_Mex.pdf}
}
\caption[5pt]{Impulse Responses to MP, CBI, and RN shocks Disentangled from FG Surprises}
\end{center}
\begin{footnotesize} 
Notes: Responses  are  presented  for a 72-month horizon with the associated 68\% highest posterior density intervals (HPDIs). FG and target-rate shocks are normalized to imply a one percentage point increase in the one-year government bond yield. The CBI and RN shocks are scaled to match the magnitude of the normalized pure FG shock in standard deviation units.
\end{footnotesize} 
\end{figure}
 

\hyperlink{Panel BF1}{Panel~B} of \hyperlink{Figure~2}{Figure~2} reports impulse responses to the different structural components disentangled from FG surprises. As before, we also plot in green the impulse responses to a contractionary pure target-rate shock. As shown in the first column, once RN and CBI components are isolated, the effects of FG shocks become more moderate and closely resemble those following target-rate shocks. This is reflected in the substantial overlap between their impulse responses. Notably, the disproportionately large and almost immediate decline in output estimated under the standard HFI approach largely disappears. In addition, inflation expectations exhibit a more muted response and appear to be better anchored after the shock. These results are more consistent with the evidence reported by \hyperlink{Solis (2023b)}{Solís (2023b)}, who shows that FG surprises play an equally important role as target-rate surprises in the transmission of monetary policy to financial markets in Mexico. They also broadly align with the international evidence reported by \hyperlink{Andrade and Ferroni (2021)}{Andrade and Ferroni (2021)} and \hyperlink{Swanson (2021, 2024)}{Swanson (2021, 2024)}, who find that FG shocks tend to generate macroeconomic effects comparable to those of conventional policy shocks.

In general, the impulse responses to CBI and RN shocks identified from FG surprises do not differ markedly from those of pure FG shocks, as shown in \hyperlink{Panel BF1}{Panel~B} of \hyperlink{Figure~2}{Figure~2}. One aspect worth highlighting, however, is the behavior of inflation expectations, which tend to increase in the short run following a favorable CBI shock and may thus help explain their tendency to rise in response to unpurged FG shocks. This contrasts with the results for conventional MP shocks and suggests that attenuation bias operating through CBI and RN channels is likely limited in the case of FG. One possible interpretation is that the central bank responds to evolving demand conditions primarily through adjustments in the policy rate rather than through communication about the future policy path. This behavior may reflect the fact that, in EMEs, FG has typically been used as a complementary rather than a primary policy tool. Unlike in AEs—where central banks were constrained by the effective lower bound and relied more heavily on this unconventional tool—FG in EMEs tends to reinforce and complement policy rate decisions, rather than substitute for them. As a result, CBI and RN shocks may play a more prominent role in attenuating the effects of target-rate surprises than those of FG surprises.

Another aspect worth highlighting is the response of domestic output to hawkish RN shocks identified from FG surprises. The dynamics closely resemble those following unpurged FG shocks, suggesting that the amplified and rapid effects of such shocks on economic activity partly reflect RN contamination. One plausible interpretation is that, in the context of FG, market participants tend to overestimate the central bank’s reaction function. For instance, when economic activity is weakening, investors may anticipate a strongly accommodative policy path, whereas the central bank—consistent with the behavior discussed above and its primary mandate of price stability—signals a relatively more hawkish stance or provides little additional guidance due to concerns about inflationary pressures. The resulting impulse responses would then reflect both the contractionary effects of an upward revision in the expected path of policy rates and underlying weakness in output that was already underway, thereby generating disproportionately large and front-loaded declines in activity.


These findings highlight the importance of disentangling FG surprises into their structural components. While the standard HFI approach suggests that FG shocks exert large and front-loaded effects, isolating RN and CBI components reveals a more moderate transmission that more closely resembles that of conventional MP shocks. In addition, inflation expectations exhibit a more muted response and appear to remain better anchored following the unconventional policy shock. Importantly, regardless of the precise mechanism behind the amplified effects of unpurged FG shocks, disentangling their structural components helps mitigate these distortions and yields estimates that appear more consistent with standard transmission mechanisms, as well as with the empirical evidence documented by related studies.


\end{document}
